% =========================================================================
% Figura (basada en «Tema 2. Figuras», diapositiva «Dipolo físico y dipolo
% puntual»): líneas de campo completas de un dipolo físico y de un dipolo
% puntual del mismo momento, como curvas de nivel de la función de corriente
% de Stokes (flujo constante entre dos líneas consecutivas). La circunferencia
% a trazos, de radio 1,5 l, marca la región fuera de la cual ambos campos son
% prácticamente indistinguibles. Datos: figuras/datos/esq_dip_*_lc.tex.
% =========================================================================
\begin{figure}[ht]
\centering
\newcommand{\esquema}[3]{% datos, eje (líneas Psi = 0), fuentes
  \begin{tikzpicture}[x=0.82cm, y=0.82cm]
    \renewcommand{\LC}[1]{\draw[magacero, line width=0.5pt] plot coordinates {##1};}%
    \fill[magfondo] (-3,-3) rectangle (3,3);
    \begin{scope}\clip (-3,-3) rectangle (3,3);
      \input{figuras/datos/#1}
      #2
    \end{scope}
    \draw[grispizarra, densely dashed, line width=0.6pt] (0,0) circle (1.5);
    \node[font=\scriptsize, text=grispizarra, fill=magfondo, fill opacity=0.85, text opacity=1, inner sep=1pt]
      at (-45:1.5) {$r=1{,}5\,\ell$};
    #3
    \draw[grispizarra!60, line width=0.8pt] (-3,-3) rectangle (3,3);
  \end{tikzpicture}}
\begin{minipage}[t]{0.47\textwidth}\centering
\esquema{esq_dip_fisico_lc.tex}%
  {\draw[magacero, line width=0.5pt] (0,0.5) -- (0,3) (0,-0.5) -- (0,-3);
   \draw[magacero, line width=0.8pt] (0,0.5) -- (0,-0.5);}%
  {\node[cargapos, minimum size=7pt, font=\tiny\bfseries] at (0,0.5) {$+$};
   \node[carganeg, minimum size=7pt, font=\tiny\bfseries] at (0,-0.5) {$-$};}
\\[2pt]{\small\scshape (a)\; dipolo físico}
\end{minipage}\hfill
\begin{minipage}[t]{0.47\textwidth}\centering
\esquema{esq_dip_puntual_lc.tex}%
  {\draw[magacero, line width=0.5pt] (0,0.05) -- (0,3) (0,-0.05) -- (0,-3);}%
  {\fill[white] (0,0) circle (0.36);
   \draw[magacero!60, line width=0.4pt] (0,0) circle (0.36);
   \draw[-{Stealth[length=3.4pt,width=3pt]}, rojoteo, line width=1.1pt] (0,-0.26) -- (0,0.28);
   \node[font=\scriptsize, text=rojoteo, inner sep=0.5pt] at (0.16,-0.08) [anchor=west] {$\vp$};}
\\[2pt]{\small\scshape (b)\; dipolo puntual}
\end{minipage}
\caption{Líneas de campo de (a)~un dipolo físico ($\pm q$ separadas $\ell$) y (b)~un dipolo puntual del mismo momento. Las líneas son las curvas
$\Psi=\text{cte}$ de la función de corriente, con $\Psi_{f}=\cos\alpha_{-}-\cos\alpha_{+}$ (ángulos medidos desde cada carga) y $\Psi_{p}=\ell\sen^{2}\theta/r$, y se trazan
con los mismos valores en ambos paneles. Dentro de la circunferencia de radio $1{,}5\,\ell$ el dipolo físico revela su estructura (las líneas nacen y mueren en las
cargas), mientras que las del dipolo puntual convergen todas en el origen (el pequeño disco blanco solo marca su posición); fuera de ella, ambas familias son prácticamente indistinguibles. Basada en~\cite{FIGURAS-T2}.}
\label{fig:esquema-dipolo}
\end{figure}
